\subsection{Local Expansion}

The local expansion is another way to represent the potential around a node.

\subsubsection{A quick look back}

We create multipole and local expansions for the nodes. A multipole expansion represents the field produced by a source node at distant points.

\[
\Phi(\mathbf{r})
=
\sum_{l=0}^{p}
\sum_{m=-l}^{+l}
M_{l}^{m}
S_{l}^{m}
\]

This represents the gravitational potential produced by the source at an observation point \(\mathbf{r}\).

Similarly, the local expansion represents the total incoming potential at an observation point \(\mathbf{r}\).

\[
\Phi(\mathbf{r})
=
\sum_{l=0}^{p}
\sum_{m=-l}^{+l}
L_{l}^{m}
R_{l}^{m}
\]

\subsubsection{How to compute the local expansion}